\chapter{The Baseline Suite}
\label{cha:baselines}

A benchmark is only as informative as the alternatives it contains. This chapter describes
the twenty-eight forecasters evaluated under the protocol of Chapter~\ref{cha:protocol}:
why each is present, what it consumes, how it was implemented, and where the
implementation departs from the published method. Chapter~\ref{cha:benchmark} reports what
they scored.

The organising principle is adversarial rather than encyclopaedic. Each model exists to
close one escape route --- one alternative explanation of a multimodal gain that a
sceptical reader would otherwise be entitled to propose. A suite assembled this way is
smaller than a survey and more useful than one, because every row has a job.

\section{Suite design}
\label{sec:suite-design}

\begin{table}[h!]
\centering
\small
\begin{tabular}{p{1.2cm} p{6.4cm} p{5.8cm}}
\hline
\textbf{Tier} & \textbf{Models} & \textbf{Question it answers} \\
\hline
T0 & Persistence, \textbf{smart persistence}, hourly climatology, seasonal naive &
Is the task trivial? What is the reference floor? \\
T1 & LightGBM (one model per quantile), TabPFN-3 & Does a tabular model already suffice? \\
T2 & MLP, DLinear, PatchTST, iTransformer, TFT & Is a simpler or a purely supervised model
better? \\
T3 & Chronos-2 zero-shot and fine-tuned (plus oracle variants), TimesFM~2.5, TiRex, TTM-R2
zero-shot and fine-tuned & Do foundation models transfer? Is any finding
backbone-specific? \\
T4 & TS-RAG, Cross-RAG, CoRA & Would retrieval or a covariate adapter close the gap without
vision? \\
T5 & Time-VLM, Aurora, VisionTS++, UniCast & Do pseudo-images or prompting suffice? \\
T6 & Solar-VLM, CrossViVit, SUNSET & Does the domain state of the art already win? \\
\hline
\end{tabular}
\caption{The baseline suite and the question each tier answers.}
\label{tab:suite}
\end{table}

The tiers span the three adaptation strategies of Section~\ref{sec:adaptation} --- full
supervision (T1--T2), zero-shot transfer (T3), and frozen-backbone adaptation (T4) --- which
is what allows Chapter~\ref{cha:benchmark} to compare them directly rather than by
citation. They also span the endogenous/exogenous distinction of
Section~\ref{sec:multimodal}: T5 contains models whose second modality is synthesised from
the numerical series, T6 models that consume genuine satellite frames. Holding the protocol
fixed across that boundary is what makes the comparison of the two classes meaningful, and
it is the comparison that produces the sharpest result in this thesis.

\paragraph{The rebuttal matrix.} Table~\ref{tab:rebuttal} states the mapping from objection
to evidence explicitly. It is reproduced here because it is the actual specification the
suite was built against.

\begin{table}[h!]
\centering
\small
\begin{tabular}{p{6.2cm} p{7.4cm}}
\hline
\textbf{Objection} & \textbf{Baseline that answers it} \\
\hline
``A simpler model would do as well'' & DLinear, MLP, persistence family \\
``The gain is just from covariates'' & CoRA, TFT, LightGBM (all covariate-rich, no vision) \\
``Retrieval would give you the same thing'' & TS-RAG, Cross-RAG \\
``The result is specific to your backbone'' & TimesFM 2.5, TiRex, TTM-R2 (three other
families) \\
``You never showed the model reads the images'' & Shuffled-frames and mismatched-plant
controls (Chapter~\ref{cha:expdesign}) \\
``Domain-specific models already do this'' & Solar-VLM, CrossViVit, SUNSET \\
``Fusion depth does not matter'' & Late-fusion arm; UniCast as published shallow prompting \\
``Pseudo-images would be cheaper'' & Time-VLM, VisionTS++ \\
\hline
\end{tabular}
\caption{Objection-to-evidence mapping. Every cell must be occupied before the multimodal
claim can be defended.}
\label{tab:rebuttal}
\end{table}

\section{Tier 0 --- reference models}
\label{sec:tier0}

Reference models establish the floor and calibrate the reader's sense of how much of the
signal is trivially predictable. On solar data this matters more than in most domains,
because the diurnal cycle alone accounts for a very large fraction of the variance
(Figure~\ref{fig:diurnal}), and a metric that rewards predicting the sun's position is
uninformative about forecasting skill.

\paragraph{Persistence.} The last observed value is carried flat across the horizon. This is
the absolute floor. At very short lead times it is surprisingly hard to beat --- power
changes slowly under stable conditions --- and its failure mode is exactly the diurnal
ramp: at 08:00 it predicts that 08:00 output will persist until 14:00.

\paragraph{Smart persistence.} The clearness index --- the ratio of observed power to a
clear-sky model --- is carried forward instead of the power itself, and the forecast is
reconstructed by multiplying that index by the clear-sky curve over the horizon. This
respects the deterministic solar geometry and assumes only that \emph{cloud conditions}
persist. It is dramatically stronger than naive persistence and is the community's standard
reference; it is the denominator of the skill score in Section~\ref{sec:metrics}. It is also
the only model in the study permitted to use clear-sky physics, because it is the physics
reference.

\paragraph{Hourly climatology.} The mean normalised power for each month, hour and minute
bin, computed on training plants. It carries no information about the current day at all,
so it detects two pathologies: a model that fails to beat it is not using its inputs, and a
suspiciously strong climatology result indicates that the evaluation windows are dominated
by average conditions rather than by the transitions that matter.

\paragraph{Seasonal naive.} The value observed at the same clock time on the previous day.
This is the standard reference in general-purpose forecasting benchmarks and is included so
that results here can be positioned against that literature.

\section{Tier 1 --- classical machine learning}
\label{sec:tier1}

\paragraph{LightGBM.} Gradient-boosted decision trees~\cite{lightgbm} over a flattened
feature vector of lagged power and covariates. A separate model is fitted per quantile with
the pinball objective, which yields a genuine probabilistic forecast rather than a point
estimate with error bars bolted on. Gradient boosting is the honest strong baseline for
tabular problems, and a deep model that cannot beat it has not justified its cost. In the
cross-plant setting it has a specific advantage: trees are insensitive to feature scaling
and therefore transfer across plants of different capacity without any special handling.

\paragraph{TabPFN-3.} A tabular foundation model~\cite{tabpfn3} in its time-series mode:
prior-fitted, performing inference in a single forward pass with no gradient-based fitting
on the target task. It is the tabular counterpart to the time-series foundation models of
Tier~3, and it tests whether the ``foundation model'' advantage, if any, is specific to
sequence architectures or is a property of large-scale prior fitting in general.

\section{Tier 2 --- supervised deep learning}
\label{sec:tier2}

All Tier-2 models are trained from scratch on the training plants and must generalise to
disjoint test plants. This is the strongest form of the ``do you actually need pretraining''
question: these models see the same data as the fine-tuned foundation models and have no
pretraining at all.

\paragraph{MLP.} A multilayer perceptron over the flattened history window and covariates.
The simplest learned baseline, present to bound how much of the achievable skill requires
any temporal structure at all.

\paragraph{DLinear.} A single linear layer applied to a trend/seasonal
decomposition~\cite{dlinear}. This is the model that embarrassed a generation of
forecasting transformers, and its presence is non-negotiable: reviewers ask for it, and on
strongly periodic signals it is frequently competitive. It consumes the target only, no
covariates, which makes it the cleanest measure of how much information is in the power
history alone.

\paragraph{PatchTST.} A channel-independent patch transformer~\cite{patchtst}: the series is
divided into patches that act as tokens, and channels are processed independently with
shared weights. It is the strongest conventional supervised patch transformer and is
architecturally the closest supervised analogue to the Chronos-2 backbone, which makes the
comparison between them informative about pretraining specifically rather than about
architecture.

\paragraph{iTransformer.} Attention is applied across \emph{variates} rather than across
time: each channel becomes a token and the attention matrix models inter-variable
structure~\cite{itransformer}. On a problem where the target is driven by a handful of
strongly correlated exogenous quantities --- irradiance, cloud cover, temperature
(Figure~\ref{fig:correlations}) --- this is an unusually well-matched inductive bias, and
Chapter~\ref{cha:benchmark} shows the consequence.

\paragraph{TFT.} The temporal fusion transformer~\cite{tft} combines gated residual
processing, variable selection and a quantile output head. It is included because it is
quantile-native, so it supplies a probabilistic comparison point outside the foundation-model
tier. The implementation here is a reduced variant, described in
Section~\ref{sec:implementation}.

\section{Tier 3 --- time-series foundation models}
\label{sec:tier3}

\paragraph{Chronos-2, zero-shot.} The backbone of the proposed architecture, evaluated with
no adaptation~\cite{chronos2}. This row is the anchor for hypothesis H0 and the honest
starting point of the whole thesis: whatever MMTSFM achieves must be measured against what
its own backbone does untouched.

\paragraph{Chronos-2, fine-tuned.} The same model fine-tuned on the training plants. This is
the single most important comparison row for the proposed architecture, because it holds the
backbone, the data and the protocol fixed and varies only the adaptation strategy. A
multimodal architecture that does not beat its own fine-tuned backbone has not demonstrated
anything about multimodality.

\paragraph{Chronos-2 oracle variants.} Two additional rows in which the model receives
\emph{observed} future weather over the forecast horizon rather than only the deterministic
covariates known in advance. These are not competitors and are labelled as such throughout;
they quantify the headroom available from perfect conditioning, which turns out to be the
largest single measurement in the study (Section~\ref{sec:bar}).

\paragraph{TimesFM 2.5.} A decoder-only foundation model from an independent
lineage~\cite{timesfm}, present so that any conclusion about zero-shot transfer is not an
artefact of one pretraining corpus.

\paragraph{TiRex.} Built on the xLSTM family rather than on attention~\cite{tirex}, and among
the strongest zero-shot performers on general-purpose leaderboards. It is the third distinct
architecture family, which is the threshold at which a negative zero-shot result becomes a
statement about the task rather than about a model.

\paragraph{TTM-R2, zero-shot and fine-tuned.} A compact mixer-based
model~\cite{ttm} of a few million parameters. It probes the scale axis directly: if
zero-shot performance is a function of capacity, a tiny model should fail and a large one
should not. The fine-tuned variant additionally shows how much of a weak zero-shot result
adaptation can recover.

\section{Tier 4 --- frozen-backbone adaptation}
\label{sec:tier4}

This tier contains the methods that compete with the thesis directly. Each keeps a
pretrained backbone frozen and adds an external mechanism for supplying task-specific
information --- exactly the strategy of the proposed architecture, differing only in what is
supplied and how deeply it is coupled.

\paragraph{TS-RAG.} Retrieval-augmented forecasting on a frozen backbone~\cite{tsrag}:
windows similar to the current context are retrieved from a datastore and fused into the
forecast. Per Section~\ref{sec:parity} the datastore contains training-plant windows only,
and the mixing weight is tuned on validation plants. This is the ``you did not need vision,
you needed precedent'' objection in runnable form.

\paragraph{Cross-RAG.} A stronger retrieval variant~\cite{crossrag} with clear-sky-aware
retrieval keys and per-step mixing weights, retrieving across series rather than only within
one. It is included so that the retrieval tier is represented by a tuned modern method
rather than by a strawman.

\paragraph{CoRA.} Covariate-aware adaptation of a frozen backbone through a zero-initialised
residual adapter~\cite{cora}. This is the closest published competitor to the proposed
architecture: same frozen-backbone premise, same idea of injecting external information
through a learned interface, but coupling the modalities at a single insertion point rather
than throughout the temporal operator. If deep token-level fusion cannot beat CoRA, the
fusion-depth claim does not survive.

\section{Tier 5 --- endogenous multimodal}
\label{sec:tier5}

Models whose second modality is synthesised from the numerical series or its covariates. No
external sensor is consumed. Their presence is what allows the thesis to separate visual
\emph{inductive bias} from visual \emph{information}.

\paragraph{Time-VLM.} The series is rendered as an image and passed, with a textual
description, to pretrained vision-language backbones~\cite{timevlm}. Nothing enters the
model that was not already in the numerical track; what is exploited is the prior structure
of a visual encoder --- locality, periodicity, multi-scale texture --- applied to a
re-rendering of the input.

\paragraph{VisionTS++.} A related approach mapping series to images to exploit a frozen
masked autoencoder backbone~\cite{visiontspp}, included to show that the endogenous strategy
is not a single lucky implementation.

\paragraph{Aurora.} A multimodal foundation model evaluated zero-shot~\cite{aurora}, with
weather text templated from the covariate track. It answers the question ``why not simply use
an off-the-shelf multimodal foundation model''.

\paragraph{UniCast.} Vision and text embeddings injected as soft prompts into a frozen
time-series foundation model~\cite{unicast}. Architecturally this is the shallowest point on
the fusion-depth axis while still being a genuine multimodal method, which makes it the
natural contrast for the interleaving hypothesis. It is placed in this tier because in the
configuration run here no external text encoder is available and the visual path operates on
generic features.

\section{Tier 6 --- exogenous multimodal, photovoltaic domain}
\label{sec:tier6}

Models that consume genuine satellite imagery, all of them from the photovoltaic
forecasting literature.

\paragraph{Solar-VLM.} The primary domain state of the art~\cite{solarvlm}: satellite
imagery, covariates and generated weather text combined through a vision-language backbone.
Per the parity rules its text is generated only from covariates available to every other
model; no external weather service is consulted.

\paragraph{CrossViVit.} The reference deep satellite-plus-series architecture, cross-attending
video tokens into a temporal transformer~\cite{crossvivit}. It is the strongest non-foundation
multimodal competitor and the closest published relative of the proposed fusion mechanism ---
both couple modalities inside the encoder rather than at a readout, differing in whether the
coupling is cross-attention or token interleaving.

\paragraph{SUNSET.} The canonical convolutional solar baseline~\cite{sunset}, a hybrid of an
image branch and a power branch. Nearly every paper in the domain compares against it, so its
presence anchors this benchmark to that literature.

\section{What was deliberately excluded}
\label{sec:exclusions}

A benchmark's exclusions are as much a design decision as its inclusions, and leaving them
unstated invites the suspicion that they were chosen after the results were known. The
following were considered and left out, each for a stated reason.

\paragraph{SPIRIT and PV-VLM.} Both are photovoltaic-domain multimodal methods that would
have strengthened Tier~6. They were excluded for effort budget rather than on principle:
Tier~6 already contains three domain models spanning convolutional, cross-attention and
vision-language designs, and a fourth would have added breadth without opening a new axis.
They remain the first additions a follow-up should make.

\paragraph{FusionSF and M3S-Net.} Non-foundation photovoltaic fusion architectures. Both are
released against corpora structured for different tasks (hourly day-ahead windows), and
porting them to this protocol would have required changes substantial enough that the
resulting row would measure the port rather than the method.

\paragraph{MEMTS and parametric-memory adapters.} These occupy the same niche as the
retrieval baselines --- external information supplied to a frozen backbone --- which is
already represented by two tuned methods. The marginal information from a third mechanism in
the same family was judged low.

\paragraph{Additional foundation-model families.} Toto and Moirai were available and were not
run. The methodological requirement was three distinct architectural families evaluated
zero-shot, which Chronos-2, TimesFM, TiRex and TTM already satisfy; a fifth transformer would
not have changed the conclusion of Section~\ref{sec:finding2}.

\paragraph{Few-shot in-context adaptation curves.} Accuracy as a function of $K$ support days
per test plant is a different protocol --- context matching rather than zero-shot transfer ---
and mixing the two in one table would obscure both. The harness supports it, and it is
available as an appendix experiment on request, with $K \in \{0,1,3,7\}$ days used either as a
retrieval datastore or as a linear-probe update.

\paragraph{Post-hoc recalibration.} Several models exhibit systematic scale bias
(Section~\ref{sec:qualitative}) that an affine correction would substantially repair. It was
not applied, because the correction would have to be fitted somewhere --- on validation plants
at best --- and the resulting table would compare forecasters-plus-post-processing rather than
forecasters. The bias is reported instead.

\section{Implementation and deviations}
\label{sec:implementation}

Third-party models were vendored and adapted to the tensor contract of
Section~\ref{sec:contract} rather than reimplemented from their papers. The deviations
below are recorded because they bound how strongly the corresponding rows may be read; a
benchmark that hides its integration compromises is not reproducible.

\begin{itemize}
  \item \textbf{Tier 2} is served by a single port of a common deep time-series
    library~\cite{tslib}, so PatchTST, iTransformer, DLinear and MLP share one trainer and
    differ only in architecture. This is a deliberate fairness choice: differences between
    those four rows cannot be attributed to training-loop details.
  \item \textbf{TFT} is a reduced variant retaining the quantile head and gated residual
    processing but not the full variable-selection stack. Its row should be read as a
    quantile-native supervised reference rather than as a faithful reproduction.
  \item \textbf{TTM} uses the R2 checkpoint, not R3. The R3 release is a trend-plus-residual
    decomposition model whose parameter keys do not load into the available library version;
    loading it silently produces a randomly initialised network that trains and evaluates
    without error. Histories shorter than the model's context are edge-padded with an
    observed mask rather than zero-padded, since zero is a meaningful power value; fine-tuned
    variants receive a frequency token.
  \item \textbf{CrossViVit} runs with single-channel imagery and synthesised coordinate
    encodings, because the \texttt{uk\_pv} product is greyscale and the original model expects
    multi-channel input with precise geospatial coordinates. Its row is a conservative
    estimate of the published architecture.
  \item \textbf{Time-VLM} runs on evaluation windows not bit-aligned with the rest of the
    suite. The number is reported and flagged rather than omitted; an aligned re-run is
    required before it is used as a precise comparison point.
  \item \textbf{Retrieval baselines} populate datastores from training plants only, with
    mixing weights tuned on validation plants.
  \item \textbf{Foundation models with fixed maximum context} consume their most recent
    supported window rather than the full fourteen days. This is stated rather than silently
    tolerated, because context length is a fairness variable
    (Section~\ref{sec:fairness}).
\end{itemize}

\section{On tuning budgets}
\label{sec:tuning}

The hardest fairness question in a benchmark of this kind is not which inputs a model
receives --- that is settled by the parity matrix --- but how much effort went into making
each one work. A model can lose because its architecture is unsuited to the task or because
nobody tuned it, and no protocol fully separates the two.

Three practices bound the problem here without solving it. Within Tier~2, four models share a
single trainer and a single hyperparameter search procedure, so differences \emph{within} that
tier are architectural rather than effort-driven --- which is what makes the iTransformer
against PatchTST comparison in Chapter~\ref{cha:benchmark} interpretable. Across tiers,
vendored models retain their authors' published defaults wherever those defaults are
applicable, on the principle that a method's own authors are better placed to configure it
than a third party with a different dataset. And the methods most threatening to the thesis
--- the retrieval baselines and the covariate adapter --- were tuned rather than accepted at
defaults, because a benchmark whose strongest competitors are strawmen proves nothing.

The residual risk runs in a specific direction, and it should be stated: the proposed
architecture received far more attention from its author than any baseline did. Where that
asymmetry would matter most --- the comparison against the same backbone fine-tuned, and
against the vision-free arm of the same model --- the comparisons are internal, sharing the
recipe, the schedule and the data, which is precisely why those two rows and not an external
baseline are treated as the decisive ones in Chapter~\ref{cha:expdesign}.

\paragraph{A correction that changed a conclusion.} One row in an earlier revision of this
benchmark reported a catastrophic score for Time-VLM. The cause was a stale artefact: the
prediction file had been re-saved after inverse scaling, but the metric file was never
regenerated, so the table showed an error computed on standardised rather than physical
units. Re-scoring from the corrected predictions --- with no retraining --- moved the model
from last place to second. All other externally integrated models re-score bit-identically to
their recorded rows, so the defect was isolated; but the episode is recorded here because it
is the reason every row in Chapter~\ref{cha:benchmark} was re-derived from saved predictions
rather than trusted from its original metric file.

\par

